% 00_abstract.tex — Abstract. Drafted by item P2-E8-1.
% Body only: this file is \input{} inside \abstract{...} in main.tex (no \section).
% NMI rule (document_check.md C): UNREFERENCED abstract — no \cite — and ~150-200 words.
% Maier/NMI voice per ../STYLE.md; honesty contract per ../plan.md.
We ask whether a faithful explanation is already an understanding.
Interpretability methods are judged by their faithfulness --- whether they name
the true causes --- yet on real neural systems there is no ground truth against
which to check them. We remove that excuse on a system whose ground truth is
computable: a bit-exact, differentiable port of the Atari~2600, never built to be
interpretable, in which the cause of any output follows from exact intervention.
Scoring about~31 interpretability methods against this oracle, we find causal
methods faithful, while popular gradient and correlational saliency methods are
provably zero on discrete position outputs (activation patching $1.000$ versus
vanilla saliency $0.000$) yet score higher on human plausibility --- a danger
zone of conviction without correctness. Yet
faithfulness is the smaller half of understanding: an exact method returns a
causal-effect map or a data-flow graph, not the computation. A graph-correct
circuit need not reproduce behavior (sufficiency $0.44$), and a feature matching a
named variable can be causally inert (matched fraction $1.0$,
faithfulness $0.04$). We name this measured
residual the \emph{semantic gap}, and stress that we \emph{measure} it, not
\emph{close} it. As such, the approach offers a reusable ground-truth benchmark
for explainable AI, mechanistic interpretability, AI safety, and software
verification. We assume that it will help redirect the field toward recovering
meaning, not merely confirming causes.

% BIBTEX-NEEDED: (none) — the abstract is intentionally citation-free per the NMI
%   unreferenced-abstract rule (document_check.md C). No \cite keys introduced;
%   references.bib unchanged.
